% ===================================================================
%  اللوحة رقم ١١ — المدينة ومبانيها. — الحريق.
%  Traduction 2026 d'apres texto/io, controlee sur texto/fr, texto/en
%  et texto/es. Consigne : voir texto/ar/10-lawha-01.tex.
% ===================================================================

%%K t11-apar-1 apar
\VUsaut{2.0mm}
\VUcentre{13.2pt}{90}{السلسلة الثالثة}
\VUsaut{3.0mm}
\VUfilet{16mm}
\VUsaut{5.6mm}
\VUtitre{15.0pt}{60}{اللوحة رقم ١١}
\VUsaut{1.4mm}
\VUpk{11.4pt}{40}{المدينة ومبانيها. --- الحريق.}
\VUsaut{2.2mm}

%%K t11-tit-1 sub
\VUpk{10.2pt}{40}{الضواحي.}
\VUsaut{3.0mm}

%%K t11-c1-01-1 p
1. --- أسكن مدينة كبيرة على بعد مئة كيلومتر من المحيط، وهي مع ذلك
\VUgras{ميناء} \textsuperscript{(1)} من الطراز الأول. والنهر الذي يجري بمحاذاتها
عرضه أربعمئة متر ويؤلف مرسى جميلًا جدًا.

%%K t11-c1-02-1 p
2. --- والجزء الممتد على \VUgras{الأرصفة} \textsuperscript{(2)} من المدينة كله له
مظهر بحري صناعي بحت، ويؤلف \VUgras{ضاحية} \textsuperscript{(3)} لا يسكنها إلا
البحّارة والعمال. وهؤلاء يعملون في \VUgras{المصانع} \textsuperscript{(4)} وفي
\VUgras{مصنع الغاز} \textsuperscript{(5)}، الذي تُرى \VUgras{مِدخنته}
\textsuperscript{(6)} العالية و\VUgras{خزان غازه} من بعيد.

%%K t11-c1-03-1 p
3. --- وفي القرون الوسطى كانت المدينة محصَّنة، ولا تزال تُوجد بعض بقايا
أسوارها، ولا سيما \VUgras{باب} \textsuperscript{(8)} \VUgras{القصر المحصَّن}
\textsuperscript{(7)} القديم يكتنفه \VUgras{برجاه} \textsuperscript{(9)}، وهما اليوم
\VUgras{سجن}.

%%K t11-c1-04-1 p
4. --- وفي هذا \VUgras{الحيّ} كله، وهو يبدو قديمًا جدًا، مبنى واحد حديث
نسبيًا: \VUgras{دار الجمرك} \textsuperscript{(10)}. وهي بين الرصيف
و\VUgras{الشارع} \textsuperscript{(11)} الصغير المؤدي إلى \VUgras{دار البلدية}
\textsuperscript{(12)} القديمة. وهذا المبنى الأخير يسهل التعرف إليه
بـ\VUgras{برج الجرس} \textsuperscript{(13)} العملاق الذي يعلو المدينة العتيقة.

%%K t11-c1-05-1 p
5. --- وعلى الجانب الآخر من الشارع \VUgras{يقوم} مبنى على طراز دار الجمرك
نفسه: إنه \VUgras{البورصة} \textsuperscript{(14)}. وإحدى واجهاته تطل على ساحة
صغيرة فيها \VUgras{دار المحافظة} \textsuperscript{(15)}. فمدينتنا، وإن لم تكن
عاصمة، مركز إقليمي مهم على كل حال.

%%K t11-tit-2 sub
\VUsaut{5.0mm}
\VUpk{10.2pt}{40}{الجزء الأعلى من المدينة.}
\VUsaut{3.4mm}

%%K t11-c1-06-1 p
6. --- والحامية، وهي كبيرة إلى حد ما، تسكن \VUgras{ثكنات} \textsuperscript{(16)}
واسعة صحية جدًا؛ وهي تمتد إلى سياج \VUgras{الحديقة البلدية} \textsuperscript{(17)}.
و\VUgras{الجادة} \textsuperscript{(19)} المؤدية إلى هذه الحديقة تمر خلف
\VUgras{صندوق التوفير} \textsuperscript{(18)} وتخرج على ساحة \VUgras{الكاتدرائية}
\textsuperscript{(20)}.

%%K t11-c1-07-1 p
7. --- وهذه المباني كلها تتدرّج على تلّ صغير تقوم عليه أيضًا
\VUgras{ثانويتنا} \textsuperscript{(21)} الكبيرة.

%%K t11-c1-07-2 p
وإذا نزلت طريقًا قليل الانحدار يلتقي بالشارع الكبير، بلغت \VUgras{قصر
العدل} \textsuperscript{(22)}، وأمامه \VUgras{حديقة عامة} \textsuperscript{(26)} لطيفة.
وهذه \VUgras{الساحة} ليست واسعة جدًا، لكنها في وسط المدينة واحة خضرة
وبرودة. ولذلك يكثر ارتيادها من أهل المدينة الذين يحبون أن يجلسوا تحت
مظلة \VUgras{المقهى} \textsuperscript{(25)} ليسمعوا الفرقة العسكرية التي تعزف
أيام الخميس والأحد في \VUgras{الجوسق} \textsuperscript{(27)} القائم بين المروج
وأحواض الزهر.

%%K t11-c1-08-1 p
8. --- وعلى أحد تلك المروج ينتصب \VUgras{تمثال} \textsuperscript{(28)} برونزي،
نصب أُقيم إحياءً لذكرى أحد أبناء مدينتنا الذين أسدوا إلى مسقط رأسهم
خدمات جليلة.

%%K t11-tit-3 sub
\VUsaut{4.0mm}
\VUpk{10.2pt}{40}{التنقل.}
\VUsaut{3.4mm}

%%K t11-c1-09-1 p
9. --- ومدينتنا لم تُضَأ بالكهرباء بعد؛ ليس فيها إلا \VUgras{مصابيح الغاز}
\textsuperscript{(29)}. لكن \VUgras{الصحافة المحلية} تحمل حملة ليُحسَّن هذا
النظام في الإنارة، \VUgras{كما} حُسِّن من قبلُ \VUgras{نظام الجرّ} في
الترام. فـ\VUgras{تلك} العربات القديمة \VUgras{التي كانت تجرّها}
\VUgras{الخيل} استُبدل بها \VUgras{ترام كهربائي} \textsuperscript{(31)} عملي ينقل
الركاب سريعًا من طرف \VUgras{المدينة إلى طرفها الآخر} بأجرة
\VUgras{زهيدة} جدًا.

%%K t11-c1-10-1 p
10. --- وقرب قصر العدل \VUgras{المصرف} \textsuperscript{(32)}، حيث تُخصم
الكمبيالات التجارية. و\VUgras{سعاة المصرف} \textsuperscript{(33)} يطوفون بالبيوت
لتحصيل ما هو مستحق.

%%K t11-tit-4 sub
\VUsaut{4.0mm}
\VUpk{10.2pt}{40}{الفن والعلم.}
\VUsaut{3.4mm}

%%K t11-c1-11-1 p
11. --- والمبنى الجميل القريب هو \VUgras{المتحف}. وهو يضم في آن واحد
\VUgras{مكتبة} \textsuperscript{(34)} المدينة، و\VUgras{مجموعات تاريخ طبيعي}
\textsuperscript{(35)} جميلة (متحف تاريخ طبيعي)، و\VUgras{صالة لوحات}
\textsuperscript{(36)} رشيقة تؤلف \VUgras{معرضًا} دائمًا رائعًا.

%%K t11-c1-11-2 p
وهو يُفتح للجمهور مرتين في الأسبوع، لكن الزائرين يستطيعون رؤيته في أي
يوم لقاء إكرامية صغيرة لـ\VUgras{الحارس} \textsuperscript{(37)}. ويأتي
\VUgras{الرسّامون} \textsuperscript{(38)} ليعملوا فيه. تراهم \VUgras{أمام}
حواملهم، و\VUgras{لوح الألوان} في أيديهم، ينسخون اللوحات المشهورة.

%%K t11-c1-12-1 p
12. --- وعلى المرتفعات خلف المتحف تلمح \VUgras{قبة} \textsuperscript{(40)}
\VUgras{المستشفى} \textsuperscript{(39)} ومباني \VUgras{دار المعلمين}
\textsuperscript{(41)} حيث تخرّج معلمو مدارسنا البلدية. وفي ساحة المسرح
\VUgras{مكتب بريد} \textsuperscript{(43)} في \VUgras{بيت خاص} \textsuperscript{(42)}.

%%K t11-tit-5 sub
\VUsaut{5.0mm}
\VUpk{11.4pt}{40}{الحريق.}
\VUsaut{3.0mm}

%%K t11-tit-6 sub
\VUpk{10.2pt}{40}{صيحة الحريق.}
\VUsaut{3.4mm}

%%K t11-c2-01-1 p
1. --- خبر مقلق يسري في المدينة: شبّ حريق في المسرح لتوّه! ولما بلغتُ
\VUgras{الساحة} \textsuperscript{(96)} كان الحشد قد تجمّع على \VUgras{الرصيف}
\textsuperscript{(46)} وعلى \VUgras{الطريق} \textsuperscript{(47)}. و\VUgras{موظفو
البريد} \textsuperscript{(44)} و\VUgras{سعاة البريد} \textsuperscript{(45)} يخرجون من
مكتب البريد. واضطر \VUgras{سائق الترام} \textsuperscript{(48)} إلى إيقاف عربته
ليفسح لـ\VUgras{سيارة إسعاف} \textsuperscript{(50)} بلدية قادمة و\VUgras{جرّاح}
\textsuperscript{(52)} و\VUgras{ممرضة} \textsuperscript{(51)} فيها، لإسعاف
\VUgras{مصاب} \textsuperscript{(53)} حُمل على \VUgras{نقّالة} \textsuperscript{(54)} إلى
موضع قرب \VUgras{كشك الجرائد} \textsuperscript{(55)}.

%%K t11-c2-02-1 p
2. --- وسريّة مشاة عائدة من التدريب توقفت لتعين على حفظ النظام مع
\VUgras{حرس البلدية} \textsuperscript{(61)} الفرسان. و\VUgras{الفرسان}، وخيلهم
تنتصب، يدفعون \VUgras{المتفرجين} \textsuperscript{(56)} أحسن مما يستطيع
\VUgras{الجنود} \textsuperscript{(57)} أو \VUgras{الدرك} \textsuperscript{(63)}. والدرك
فوق ذلك مشغولون جدًا. فأحدهم يدلّ \VUgras{رجال الشرطة} \textsuperscript{(64)} على
المكان الذي يُحمل إليه \VUgras{مصاب} \textsuperscript{(65)} آخر، وهو إطفائي شجاع
سقط عن سلّم.

%%K t11-c2-03-1 p
3. --- و\VUgras{الضابط} \textsuperscript{(66)} يبيّن بالضبط أين توضع
\VUgras{المضخة البخارية} \textsuperscript{(67)} القادمة الآن. وفي انتظار تلك
الآلة القوية أمر بنصب \VUgras{مضخة يدوية} \textsuperscript{(68)}
\VUgras{خرطومها} \textsuperscript{(69)} الطويل يصبّ على النار سيولًا من الماء.
ويشتغل عليها ستة إطفائيين، بينما يمسك رفيقهم \VUgras{الفوهة}
\textsuperscript{(70)} ويصوّب \VUgras{النفّاث} \textsuperscript{(71)} إلى قلب الحريق.

%%K t11-c2-04-1 p
4. --- وعلى الجانب الآخر من المسرح رُفع \VUgras{السلّم} \textsuperscript{(72)}
الكبير. وهو يتيح للإطفائيين الاقتراب من الألسنة المتصاعدة بين دوّامات
\VUgras{الدخان} \textsuperscript{(75)} الأسود.

%%K t11-tit-7 sub
\VUsaut{5.0mm}
\VUpk{10.2pt}{40}{أعمال شجاعة.}
\VUsaut{3.4mm}

%%K t11-c2-05-1 p
5. --- شبّت \VUgras{النار} \textsuperscript{(74)} في بهو \VUgras{المسرح}
\textsuperscript{(73)} أثناء تمرين على \VUgras{الأوبرا} التي كان \VUgras{عرضها}
الأول سيجري غدًا. ولحسن الحظ لم يكن في القاعة إلا بضعة \VUgras{متفرجين}
\textsuperscript{(76)}. وقد لجأ هؤلاء المساكين إلى \VUgras{الشرفة} \textsuperscript{(77)}،
حيث يأتي \VUgras{إطفائيون} \textsuperscript{(86)} شجعان لنجدتهم بسلّم وحبال.

%%K t11-c2-06-1 p
6. --- وتحت \VUgras{الرواق} \textsuperscript{(78)}، حيث يعلن \VUgras{الملصق}
\textsuperscript{(79)} عن العرض، يُرى \VUgras{عازفو} \textsuperscript{(81)}
\VUgras{الأوركسترا} \textsuperscript{(80)} يفرّون حاملين آلاتهم:
\VUgras{الكمان} \textsuperscript{(81)}، و\VUgras{الناي} \textsuperscript{(82)}،
و\VUgras{التشيلو} \textsuperscript{(83)}، و\VUgras{الترومبون} \textsuperscript{(84)}،
و\VUgras{الطبل} \textsuperscript{(85)}، وهكذا. ويُبذل جهد لإنقاذ جزء من
\VUgras{الأثاث} \textsuperscript{(87)}.

%%K t11-c2-07-1 p
7. --- و\VUgras{الممثلون} \textsuperscript{(89)} و\VUgras{الممثلات}
\textsuperscript{(88)}، و\VUgras{المغنّون} والمغنّيات، اضطروا إلى الخروج
مسرعين في \VUgras{أزيائهم} \textsuperscript{(90)} المسرحية. والتينور يشرح
لـ\VUgras{صحفيّ} \textsuperscript{(92)} كيف وقع الأمر. و\VUgras{المُراسِل} جاء
راكضًا من اللحظة الأولى مع رجال السلطة: \VUgras{المحافظ} \textsuperscript{(91)}،
و\VUgras{العمدة} \textsuperscript{(93)}، و\VUgras{مدير الشرطة} \textsuperscript{(94)}،
و\VUgras{قاضٍ} \textsuperscript{(95)} سيجري تحقيقًا لتحديد أين تقع المسؤولية.
\VUsaut{6.0mm}
\VUfilet{22mm}
